\documentclass[12pt]{article}
\usepackage{amsmath}
\usepackage{amsfonts}
\usepackage[margin=1in]{geometry}

\title{APM1220 FA1}
\author{Kristen Jelaena Cesista}

\begin{document}

\maketitle

\noindent\rule{\textwidth}{0.4pt}

\vspace{2em}

\noindent {\Large \textbf{Problem 2.2}} \hfill {\large \textbf{(5 points)}}

\vspace{1em}

\noindent Given the matrices

$$
\mathbf{A} = \begin{bmatrix} -1 & 3 \\ 4 & 2 \end{bmatrix}, \quad
\mathbf{B} = \begin{bmatrix} 4 & -3 \\ 1 & -2 \\ -2 & 0 \end{bmatrix}, \quad \text{and} \quad
\mathbf{C} = \begin{bmatrix} 5 \\ -4 \\ 2 \end{bmatrix}
$$

\noindent perform the indicated multiplications:

\begin{enumerate}
\item[(a)] $5\mathbf{A}$

$$
5A
=
5
\begin{bmatrix}
-1 & 3\\
4 & 2
\end{bmatrix}
=
\begin{bmatrix}
-1\cdot5 & 3\cdot5\\
4\cdot5 & 2\cdot5
\end{bmatrix}
$$

$$
\boxed{
5A=
\begin{bmatrix}
-5 & 15\\
20 & 10
\end{bmatrix}
}
$$

\item[(b)] $\mathbf{B}\mathbf{A}$

$$
BA
=
\begin{bmatrix}
4 & -3\\
1 & -2\\
-2 & 0
\end{bmatrix}
\begin{bmatrix}
-1 & 3\\
4 & 2
\end{bmatrix}
$$

$$
=
\begin{bmatrix}
4(-1)+(-3)(4) & 4(3)+(-3)(2)\\
1(-1)+(-2)(4) & 1(3)+(-2)(2)\\
(-2)(-1)+(0)(4) & (-2)(3)+(0)(2)
\end{bmatrix}
$$

$$
=
\begin{bmatrix}
-4-12 & 12-6\\
-1-8 & 3-4\\
2+0 & -6+0
\end{bmatrix}
$$

$$
\boxed{
BA=
\begin{bmatrix}
-16 & 6\\
-9 & -1\\
2 & -6
\end{bmatrix}
}
$$

\item[(c)] $\mathbf{A}'\mathbf{B}'$

$$
A'
=
\begin{bmatrix}
-1 & 4\\
3 & 2
\end{bmatrix}
\qquad
B'
=
\begin{bmatrix}
4 & 1 & -2\\
-3 & -2 & 0
\end{bmatrix}
$$

$$
A'B'
=
\begin{bmatrix}
-1 & 4\\
3 & 2
\end{bmatrix}
\begin{bmatrix}
4 & 1 & -2\\
-3 & -2 & 0
\end{bmatrix}
$$

$$
=
\begin{bmatrix}
(-1)(4)+(4)(-3) & (-1)(1)+(4)(-2) & (-1)(-2)+(4)(0)\\
(3)(4)+(2)(-3) & (3)(1)+(2)(-2) & (3)(-2)+(2)(0)
\end{bmatrix}
$$

$$
=
\begin{bmatrix}
-4-12 & -1-8 & 2+0\\
12-6 & 3-4 & -6+0
\end{bmatrix}
$$

$$
\boxed{
A'B'
=
\begin{bmatrix}
-16 & -9 & 2\\
6 & -1 & -6
\end{bmatrix}
}
$$

\item[(d)] $\mathbf{C}'\mathbf{B}$

$$
C'
=
\begin{bmatrix}
5 & -4 & 2
\end{bmatrix}
$$

$$
C'B
=
\begin{bmatrix}
5 & -4 & 2
\end{bmatrix}
\begin{bmatrix}
4 & -3\\
1 & -2\\
-2 & 0
\end{bmatrix}
$$

$$
=
\begin{bmatrix}
(5)(4)+(-4)(1)+(2)(-2)
&
(5)(-3)+(-4)(-2)+(2)(0)
\end{bmatrix}
$$

$$
=
\begin{bmatrix}
20-4-4 & -15+8+0
\end{bmatrix}
$$

$$
\boxed{
C'B=
\begin{bmatrix}
12 & -7
\end{bmatrix}
}
$$

\item[(e)] Is $\mathbf{A}\mathbf{B}$ defined?

\noindent For a matrix multiplication to be defined, the number of columns of the 1st matrix must equal the number of rows of the 2nd one.

\noindent The matrix $\mathbf{A}$ has dimensions $2\times2$, while the matrix $\mathbf{B}$ has dimensions $3\times2$. $\mathbf{B}\mathbf{A}$ is defined because the number of columns in $\mathbf{B}$ and the number of rows in $\mathbf{A}$ are both 2. Similarly, $\mathbf{A}'\mathbf{B}'$ is defined because the number of columns in $\mathbf{A}'$ is 2 and the number of rows in $\mathbf{B}'$ is also 2.

\noindent For the multiplication $\mathbf{A}\mathbf{B}$ to be defined, the number of columns of $\mathbf{A}$ must equal the number of rows of $\mathbf{B}$. The number of columns of $\mathbf{A}$ is $2$, while the number of rows of $\mathbf{B}$ is $3$. Since

$$
2\neq3,
$$

\noindent the multiplication is not defined.

$$
\boxed{\mathbf{A}\mathbf{B}\text{ is not defined.}}
$$

\end{enumerate}

\vspace{2em}

\noindent\rule{\textwidth}{0.4pt}

\vspace{2em}

\noindent {\Large \textbf{Problem 2.4}} \hfill {\large \textbf{(6 points)}}

\vspace{1em}

\noindent When $\mathbf{A}^{-1}$ and $\mathbf{B}^{-1}$ exist, prove each of the following:

\begin{enumerate}
\item[(a)] $(\mathbf{A}')^{-1} = (\mathbf{A}^{-1})'$

\vspace{0.5em}

\noindent Since $\mathbf{A}^{-1}$ exists, by definition of an inverse we know that

$$
\mathbf{A}\mathbf{A}^{-1} = \mathbf{I}.
$$

\noindent Taking the transpose of both sides gives

$$
(\mathbf{A}\mathbf{A}^{-1})' = \mathbf{I}'.
$$

\noindent Since the transpose of an identity matrix is itself, $\mathbf{I}' = \mathbf{I}$. Also, the transpose of a product reverses the order of the matrices, so

$$
\underline{
(\mathbf{A}^{-1})'\mathbf{A}' = \mathbf{I}
}
$$

\noindent Similarly, using the other side of the inverse definition,

$$
\mathbf{A}^{-1}\mathbf{A} = \mathbf{I},
$$

\noindent taking the transpose of both sides will result in

$$
(\mathbf{A}^{-1}\mathbf{A})' = \mathbf{I}',
$$

\noindent and applying the same properties  as before will give us

$$
\underline{
\mathbf{A}'(\mathbf{A}^{-1})' = \mathbf{I}
}
$$

\noindent We showed that $(\mathbf{A}^{-1})'$ produces the identity matrix $\mathbf{I}$ when multiplied by $\mathbf{A}'$ on both the left and the right. By definition, it is the inverse of $\mathbf{A}'$. Therefore,

$$
\boxed{
(\mathbf{A}')^{-1}=(\mathbf{A}^{-1})'
}
$$

\item[(b)] $(\mathbf{A}\mathbf{B})^{-1} = \mathbf{B}^{-1}\mathbf{A}^{-1}$

\vspace{0.5em}

\noindent Consider the product

$$
(\mathbf{B}^{-1}\mathbf{A}^{-1})(\mathbf{A}\mathbf{B}).
$$

\noindent Using associativity property of matrix multiplication, we get

$$
(\mathbf{B}^{-1}\mathbf{A}^{-1})(\mathbf{A}\mathbf{B})
=
\mathbf{B}^{-1}(\mathbf{A}^{-1}\mathbf{A})\mathbf{B}.
$$

\noindent Since $\mathbf{A}^{-1}\mathbf{A} = \mathbf{I}$,

$$
\mathbf{B}^{-1}(\mathbf{A}^{-1}\mathbf{A})\mathbf{B}
=
\mathbf{B}^{-1}\mathbf{I}\mathbf{B}.
$$

\noindent Since multiplying by the identity matrix simply results in the original matrix, we have

$$
\mathbf{B}^{-1}\mathbf{I}\mathbf{B}
=
\mathbf{B}^{-1}\mathbf{B}
$$

By definition of an inverse we know that $\mathbf{B}^{-1}\mathbf{B} = \mathbf{I}$ 

\noindent Therefore,

$$
\underline{
(\mathbf{B}^{-1}\mathbf{A}^{-1})(\mathbf{A}\mathbf{B})=\mathbf{I}
}
$$

\noindent Similarly,

$$
(\mathbf{A}\mathbf{B})(\mathbf{B}^{-1}\mathbf{A}^{-1})
=
\mathbf{A}(\mathbf{B}\mathbf{B}^{-1})\mathbf{A}^{-1}
$$

\noindent Since $\mathbf{B}^{-1}\mathbf{B} = \mathbf{I}$,

$$
(\mathbf{A}\mathbf{B})(\mathbf{B}^{-1}\mathbf{A}^{-1})
=
\mathbf{A}\mathbf{I}\mathbf{A}^{-1}
$$

\noindent Since multiplying by the identity matrix simply results in the original matrix, we have

$$
\mathbf{A}^{-1}\mathbf{I}\mathbf{A}
=
\mathbf{A}^{-1}\mathbf{A}
$$

By definition of an inverse we know that $\mathbf{A}^{-1}\mathbf{A} = \mathbf{I}$ 

\noindent Therefore,

$$
\underline{
(\mathbf{A}\mathbf{B})(\mathbf{B}^{-1}\mathbf{A}^{-1})=\mathbf{I}
}
$$

\noindent Since $\mathbf{B}^{-1}\mathbf{A}^{-1}$ produces the identity matrix $\mathbf{I}$ when multiplied by $\mathbf{A}\mathbf{B}$ on both the left and the right, it is, by definition, the inverse of $\mathbf{A}\mathbf{B}$. Therefore,

$$
\boxed{
(\mathbf{A}\mathbf{B})^{-1}=\mathbf{B}^{-1}\mathbf{A}^{-1}
}
$$

\end{enumerate}

\vspace{2em}

\noindent\rule{\textwidth}{0.4pt}

\vspace{2em}

\noindent {\Large \textbf{Problem 2.8}} \hfill {\large \textbf{(6 points)}}

\vspace{1em}

\noindent Given the matrix

$$
\mathbf{A} = \begin{bmatrix} 1 & 2 \\ 2 & -2 \end{bmatrix}
$$

\noindent find the eigenvalues $\lambda_1$ and $\lambda_2$ and the associated normalized eigenvectors $\mathbf{e}_1$ and $\mathbf{e}_2$. Determine the spectral decomposition of $\mathbf{A}$.

\vspace{1em}

\noindent \textbf{Finding the eigenvalues}

\noindent The eigenvalues of matrix $\mathbf{A}$ can be found by solving the equation: $det( \lambda \mathbf{I}-\mathbf{A})$

$$
det\left(
\lambda\begin{bmatrix} 1 & 0 \\ 0 & 1 \end{bmatrix} - \begin{bmatrix} 1 & 2 \\ 2 & -2 \end{bmatrix}
\right) = 0
$$

$$
det\left(
\begin{bmatrix} \lambda -1 & -2 \\ -2 & \lambda + 2 \end{bmatrix}
\right) = 0
$$

\noindent Computing the determinant, 

$$
(\lambda-1)(\lambda+2)-(-2)(-2)=0
$$

$$
\lambda^2+\lambda-2-4=0
$$

$$
\lambda^2+\lambda-6=0
$$

\noindent Factoring the quadratic gives

$$
(\lambda-2)(\lambda+3)=0
$$

\noindent Therefore, the eigenvalues are

$$
\boxed{\lambda_1=2,\qquad \lambda_2=-3}
$$

\vspace{1em}

\noindent To solve for the eigenvectors, the equation $(\mathbf{A}-\lambda\mathbf{I})e = 0$ will be used.

\vspace{1em}

\noindent \textbf{Finding the eigenvector for $\lambda_1 = 2$}:

\noindent We solve

$$
(\mathbf{A}-2\mathbf{I})\mathbf{e}_1=\mathbf{0}.
$$

$$
\begin{bmatrix}
1-2 & 2\\
2 & -2-2
\end{bmatrix}
\begin{bmatrix}
x\\
y
\end{bmatrix}
=
\begin{bmatrix}
0\\
0
\end{bmatrix}
$$

$$
\begin{bmatrix}
-1 & 2\\
2 & -4
\end{bmatrix}
\begin{bmatrix}
x\\
y
\end{bmatrix}
=
\begin{bmatrix}
0\\
0
\end{bmatrix}
$$

\noindent From the first row,

$$
-x+2y=0
$$

$$
x=2y.
$$

\noindent Let $y=1$. Then $x=2$, so an eigenvector is

$$
\begin{bmatrix}
2\\
1
\end{bmatrix}.
$$

\noindent Its length is

$$
\sqrt{2^2+1^2}=\sqrt{5}.
$$

\noindent Therefore, the normalized eigenvector is

$$
\boxed{
\mathbf{e}_1=
\begin{bmatrix}
2/\sqrt{5}\\
1/\sqrt{5}
\end{bmatrix}
}
$$

\vspace{1em}

\noindent \textbf{Finding the eigenvector for $\lambda_2 = -3$}:

\noindent We solve

$$
(\mathbf{A}+3\mathbf{I})\mathbf{e}_2=\mathbf{0}.
$$

$$
\begin{bmatrix}
1+3 & 2\\
2 & -2+3
\end{bmatrix}
\begin{bmatrix}
x\\
y
\end{bmatrix}
=
\begin{bmatrix}
0\\
0
\end{bmatrix}
$$

$$
\begin{bmatrix}
4 & 2\\
2 & 1
\end{bmatrix}
\begin{bmatrix}
x\\
y
\end{bmatrix}
=
\begin{bmatrix}
0\\
0
\end{bmatrix}
$$

\noindent From the first row,

$$
4x+2y=0
$$

$$
y=-2x.
$$

\noindent Let $x=1$. Then $y=-2$, so an eigenvector is

$$
\begin{bmatrix}
1\\
-2
\end{bmatrix}.
$$

\noindent Its length is

$$
\sqrt{1^2+(-2)^2}=\sqrt{5}.
$$

\noindent Therefore, the normalized eigenvector is

$$
\boxed{
\mathbf{e}_2=
\begin{bmatrix}
1/\sqrt{5}\\
-2/\sqrt{5}
\end{bmatrix}
}
$$

\vspace{1em}

\noindent \textbf{Spectral decomposition}

\noindent The spectral decomposition (2-16) is

$$
\mathbf{A} = \lambda_1 \mathbf{e}_1 \mathbf{e}_1' + \lambda_2 \mathbf{e}_2 \mathbf{e}_2'.
$$

$$
\mathbf{e}_1\mathbf{e}_1' =
\begin{bmatrix} 2/\sqrt5 \\ 1/\sqrt5 \end{bmatrix}
\begin{bmatrix} 2/\sqrt5 & 1/\sqrt5 \end{bmatrix}
=
\frac{1}{5}\begin{bmatrix} 4 & 2 \\ 2 & 1 \end{bmatrix}
$$

$$
\mathbf{e}_2\mathbf{e}_2' =
\begin{bmatrix} 1/\sqrt5 \\ -2/\sqrt5 \end{bmatrix}
\begin{bmatrix} 1/\sqrt5 & -2/\sqrt5 \end{bmatrix}
=
\frac{1}{5}\begin{bmatrix} 1 & -2 \\ -2 & 4 \end{bmatrix}
$$

$$
\mathbf{A} = 2\cdot\frac{1}{5}\begin{bmatrix} 4 & 2 \\ 2 & 1 \end{bmatrix} + (-3)\cdot\frac{1}{5}\begin{bmatrix} 1 & -2 \\ -2 & 4 \end{bmatrix}
=
\begin{bmatrix} 8/5 & 4/5 \\ 4/5 & 2/5 \end{bmatrix}
+
\begin{bmatrix} -3/5 & 6/5 \\ 6/5 & -12/5 \end{bmatrix}
$$

$$
\boxed{
\mathbf{A} = 2\left(\frac{1}{5}\begin{bmatrix} 4 & 2 \\ 2 & 1 \end{bmatrix}\right) + (-3)\left(\frac{1}{5}\begin{bmatrix} 1 & -2 \\ -2 & 4 \end{bmatrix}\right)
=
\begin{bmatrix} 1 & 2 \\ 2 & -2 \end{bmatrix}
}
$$

\noindent which recovers $\mathbf{A}$, confirming the decomposition is correct.

\vspace{2em}

\noindent\rule{\textwidth}{0.4pt}

\vspace{2em}

\noindent {\Large \textbf{Problem 2.11}} \hfill {\large \textbf{(5 points)}}

\vspace{1em}

\noindent Show that the determinant of the $p \times p$ diagonal matrix $\mathbf{A} = \{a_{ij}\}$ with $a_{ij} = 0, i \neq j$, is given by the product of the diagonal elements; thus, $|\mathbf{A}| = a_{11}a_{22}\cdots a_{pp}$.

\vspace{1em}

\noindent \textbf{Proof}

\noindent Starting with cofactor expansion along the first row, using Definition 2A.24:

$$
|\mathbf{A}| = a_{11}A_{11} + a_{12}A_{12} + \cdots + a_{1p}A_{1p}
$$

\noindent where $A_{1j}$ is the cofactor of $a_{1j}$. Since $\mathbf{A}$ is diagonal, every term in the first row is zero except the first term. Thus

$$
|\mathbf{A}| = a_{11}A_{11} + 0 + \cdots + 0
$$

$$
|\mathbf{A}| = a_{11}A_{11}
$$

\noindent $\mathbf{A}_{11}$, is also the submatrix obtained by deleting the first row and first column of $\mathbf{A}$. Because $\mathbf{A}$ is diagonal, $\mathbf{A}_{11}$ is itself a diagonal matrix, with diagonal entries $a_{22}, a_{33}, \ldots, a_{pp}$:

$$
\mathbf{A}_{11} =
\begin{bmatrix}
a_{22} & 0 & \cdots & 0 \\
0 & a_{33} & \cdots & 0 \\
\vdots & & \ddots & \vdots \\
0 & 0 & \cdots & a_{pp}
\end{bmatrix}
$$

\noindent Again if we expand $|\mathbf{A}_{11}|$ along its first row, which again has only one nonzero entry, $a_{22}$, this gives us

$$
\mathbf{A}_{11} = a_{22}\mathbf{A}_{22}
$$

Thus, 
$$
|\mathbf{A}| = a_{11}a_{22}A_{22}
$$

\noindent Continuing this process, we will obtain, after $p-1$ reductions,

$$
|\mathbf{A}| = a_{11}\cdot a_{22} \cdot \mathbf{A}_{22} = a_{11}a_{22}a_{33}\cdots a_{p-1,p-1}\cdot a_{pp}.
$$

\noindent Therefore,

$$
\boxed{|\mathbf{A}| = a_{11}a_{22}\cdots a_{pp}}
$$

\vspace{2em}

\noindent\rule{\textwidth}{0.4pt}

\vspace{2em}

\noindent {\Large \textbf{Problem 2.24}} \hfill {\large \textbf{(5 points)}}

\vspace{1em}

\noindent Let $\mathbf{X}$ have covariance matrix

$$
\boldsymbol{\Sigma} = \begin{bmatrix} 4 & 0 & 0 \\ 0 & 9 & 0 \\ 0 & 0 & 1 \end{bmatrix}
$$

\noindent Find:

\begin{enumerate}
\item[(a)] $\boldsymbol{\Sigma}^{-1}$

\vspace{0.5em}

\noindent Since $\boldsymbol{\Sigma}$ is a diagonal matrix, its inverse is also a diagonal matrix. The diagonal entries of the inverse can simply be obtained by taking the reciprocals of the entries of the original matrix, thus,

$$
\boxed{
\boldsymbol{\Sigma}^{-1} =
\begin{bmatrix} 1/4 & 0 & 0 \\ 0 & 1/9 & 0 \\ 0 & 0 & 1 \end{bmatrix}
}
$$

\item[(b)] The eigenvalues and eigenvectors of $\boldsymbol{\Sigma}$

\vspace{0.5em}

\noindent For a diagonal matrix such as $\boldsymbol{\Sigma}$, the eigenvalues $\lambda_{i}$ are simply the diagonal entries $d_{ii}$, and the eigenvectors $e_{i}$ are simply the corresponding basis vectors where the i-th component is 1 and the rest is zero. Thus,

$$
\boxed{
\lambda_1 = 4,\ \mathbf{e}_1 = \begin{bmatrix} 1 \\ 0 \\ 0 \end{bmatrix}; \qquad
\lambda_2 = 9,\ \mathbf{e}_2 = \begin{bmatrix} 0 \\ 1 \\ 0 \end{bmatrix}; \qquad
\lambda_3 = 1,\ \mathbf{e}_3 = \begin{bmatrix} 0 \\ 0 \\ 1 \end{bmatrix}
}
$$

\item[(c)] The eigenvalues and eigenvectors of $\boldsymbol{\Sigma}^{-1}$

\vspace{0.5em}

\noindent Again since $\boldsymbol{\Sigma}^{-1}$ is also a diagonal matrix, its the eigenvalues $\lambda_{i}$ are simply the diagonal entries $d_{ii}$, and the eigenvectors $e_{i}$ are simply the corresponding basis vectors where the i-th component is 1 and the rest is zero. Thus,

$$
\boxed{
\lambda_1= \frac{1}{4},\ \mathbf{e}_1 = \begin{bmatrix} 1 \\ 0 \\ 0 \end{bmatrix}; \qquad
\lambda_2 = \frac{1}{9},\ \mathbf{e}_2 = \begin{bmatrix} 0 \\ 1 \\ 0 \end{bmatrix}; \qquad
\lambda_3 = 1,\ \mathbf{e}_3 = \begin{bmatrix} 0 \\ 0 \\ 1 \end{bmatrix}
}
$$

\end{enumerate}

\vspace{2em}

\noindent\rule{\textwidth}{0.4pt}

\vspace{2em}

\noindent {\Large \textbf{Problem 2.30}} \hfill {\large \textbf{(8 points)}}

\vspace{1em}

\noindent You are given the random vector $\mathbf{X}' = [X_1, X_2, X_3, X_4]$ with mean vector $\boldsymbol{\mu}_X' = [4, 3, 2, 1]$ and variance-covariance matrix

$$
\boldsymbol{\Sigma}_X =
\begin{bmatrix}
3 & 0 & 2 & 2 \\
0 & 1 & 1 & 0 \\
2 & 1 & 9 & -2 \\
2 & 0 & -2 & 4
\end{bmatrix}
$$

\noindent Partition $\mathbf{X}$ as $\mathbf{X} = \begin{bmatrix} X_1 \\ X_2 \\ X_3 \\ X_4 \end{bmatrix} = \begin{bmatrix} \mathbf{X}^{(1)} \\ \mathbf{X}^{(2)} \end{bmatrix}$, so that

$$
\boldsymbol{\mu}^{(1)} = \begin{bmatrix} 4 \\ 3 \end{bmatrix}, \quad
\boldsymbol{\mu}^{(2)} = \begin{bmatrix} 2 \\ 1 \end{bmatrix}, \quad
\boldsymbol{\Sigma}_{11} = \begin{bmatrix} 3 & 0 \\ 0 & 1 \end{bmatrix}, \quad
\boldsymbol{\Sigma}_{12} = \begin{bmatrix} 2 & 2 \\ 1 & 0 \end{bmatrix}, \quad
\boldsymbol{\Sigma}_{22} = \begin{bmatrix} 9 & -2 \\ -2 & 4 \end{bmatrix}
$$

\noindent Let $\mathbf{A} = \begin{bmatrix} 1 & 2 \end{bmatrix}$ and $\mathbf{B} = \begin{bmatrix} 1 & -2 \\ 2 & -1 \end{bmatrix}$, and consider the linear combinations $\mathbf{A}\mathbf{X}^{(1)}$ and $\mathbf{B}\mathbf{X}^{(2)}$. Find:

\begin{enumerate}
\item[(a)] $E(\mathbf{X}^{(1)})$

\vspace{0.5em}

\noindent $E(\mathbf{X}^{(1)})$ is simply the first two entries of $\boldsymbol{\mu}_X$:

$$
\boxed{
E(\mathbf{X}^{(1)}) = \begin{bmatrix} 4 \\ 3 \end{bmatrix}
}
$$

\item[(b)] $E(\mathbf{A}\mathbf{X}^{(1)})$

\vspace{0.5em}

\noindent Using $E(\mathbf{A}\mathbf{X}) = \mathbf{A}E(\mathbf{X})$:

$$
E(\mathbf{A}\mathbf{X}^{(1)}) = \mathbf{A}E(\mathbf{X}^{(1)})
=
\begin{bmatrix} 1 & 2 \end{bmatrix}
\begin{bmatrix} 4 \\ 3 \end{bmatrix}
= (1)(4) + (2)(3) = 4 + 6
$$

$$
\boxed{
E(\mathbf{A}\mathbf{X}^{(1)}) = 10
}
$$

\item[(c)] $\text{Cov}(\mathbf{X}^{(1)})$

\vspace{0.5em}

\noindent The covariance matrix of $X^{(1)}$ is simply the upper-left $2\times2$ block of $\Sigma_X$:

$$
\boxed{
\text{Cov}(\mathbf{X}^{(1)}) = \boldsymbol{\Sigma}_{11} = \begin{bmatrix} 3 & 0 \\ 0 & 1 \end{bmatrix}
}
$$

\item[(d)] $\text{Cov}(\mathbf{A}\mathbf{X}^{(1)})$

\vspace{0.5em}

\noindent Using $\text{Cov}(\mathbf{A}\mathbf{X}) = \mathbf{A}\,\text{Cov}(\mathbf{X})\,\mathbf{A}'$:

Since 

$$
\mathbf{A} = \begin{bmatrix} 1 & 2 \end{bmatrix}, \hspace{1em} A'= \begin{bmatrix} 1 \\ 2 \end{bmatrix}
$$

we get

$$
\text{Cov}(\mathbf{AX}^{(1)}) = \begin{bmatrix} 1 & 2 \end{bmatrix} \begin{bmatrix} 3 & 0 \\ 0 & 1 \end{bmatrix} \begin{bmatrix} 1 \\ 2 \end{bmatrix}
$$

Multiplying the first two matrices,

$$
\begin{bmatrix} 1 & 2 \end{bmatrix} \begin{bmatrix} 3 & 0 \\ 0 & 1 \end{bmatrix}  = \begin{bmatrix} 3 & 2
\end{bmatrix}
$$

Then multiplying the resulting product and the last matrix, we get

$$
\begin{bmatrix} 3 & 2 \end{bmatrix} \begin{bmatrix} 1 \\ 2 \end{bmatrix} = 3 + 4 =7
$$

$$
\boxed{
\text{Cov}(\mathbf{A}\mathbf{X}^{(1)}) = 7
}
$$

\item[(e)] $E(\mathbf{X}^{(2)})$

\noindent This is simply the last two elements of the mean vector

$$
\boxed{
E(\mathbf{X}^{(2)}) = \begin{bmatrix} 2 \\ 1 \end{bmatrix}
}
$$

\item[(f)] $E(\mathbf{B}\mathbf{X}^{(2)})$

\vspace{0.5em}

\noindent Using $E(\mathbf{B}\mathbf{X}^{(2)}) = \mathbf{B}\mathbf{X}^{(2)}$

So we have,
$$
\begin{bmatrix} 1 & -2 \\ 2 & -1 \end{bmatrix}
\begin{bmatrix} 2 \\ 1 \end{bmatrix}
=
\begin{bmatrix} (1)(2)+(-2)(1) \\ (2)(2)+(-1)(1) \end{bmatrix}
=
\begin{bmatrix} 2-2 \\ 4-1 \end{bmatrix} = \begin{bmatrix} 0 \\ 3 \end{bmatrix}
$$

$$
\boxed{
E(\mathbf{B}\mathbf{X}^{(2)}) = \begin{bmatrix} 0 \\ 3 \end{bmatrix}
}
$$

\item[(g)] $\text{Cov}(\mathbf{X}^{(2)})$

\noindent This is simply the bottom right 2x2 block of the mean vector
$$
\boxed{
\text{Cov}(\mathbf{X}^{(2)}) = \boldsymbol{\Sigma}_{22} = \begin{bmatrix} 9 & -2 \\ -2 & 4 \end{bmatrix}
}
$$

\item[(h)] $\text{Cov}(\mathbf{B}\mathbf{X}^{(2)})$

\vspace{0.5em}

\noindent Using $\text{Cov}(\mathbf{B}\mathbf{X}) = \mathbf{B}\,\text{Cov}(\mathbf{X})\,\mathbf{B}'$:

We have,
$$
\mathbf{B} = \begin{bmatrix} 1 & -2 \\ 2 & -1 \end{bmatrix}, \hspace{1em} \mathbf{B'} = \begin{bmatrix} 1 & 2 \\ -2 & -1 \end{bmatrix}
$$

\noindent So, 

$$
\text{Cov}(\mathbf{B}\mathbf{X}^{(2)}) = \begin{bmatrix} 1 & -2 \\ 2 & -1 \end{bmatrix} \begin{bmatrix} 9 & -2 \\ -2 & 4 \end{bmatrix}  \begin{bmatrix} 1 & 2 \\ -2 & -1 \end{bmatrix}
$$

Multiplying the first two matrices, we get

$$
\begin{bmatrix} 1 & -2 \\ 2 & -1 \end{bmatrix}
\begin{bmatrix} 9 & -2 \\ -2 & 4 \end{bmatrix}
=
\begin{bmatrix} (1)(9)+(-2)(-2) & (1)(-2)+(-2)(4) \\ (2)(9)+(-1)(-2) & (2)(-2)+(-1)(4) \end{bmatrix}
=
\begin{bmatrix} 13 & -10 \\ 20 & -8 \end{bmatrix}
$$

Now, we multiply the result with the last matrix getting

$$
\begin{bmatrix} 13 & -10 \\ 20 & -8 \end{bmatrix}
\begin{bmatrix} 1 & 2 \\ -2 & -1 \end{bmatrix}
=
\begin{bmatrix} (13)(1)+(-10)(-2) & (13)(2)+(-10)(-1) \\ (20)(1)+(-8)(-2) & (20)(2)+(-8)(-1) \end{bmatrix} = \begin{bmatrix} 33 & 36 \\ 36 & 48 \end{bmatrix}
$$

Therefore,

$$
\boxed{
\text{Cov}(\mathbf{B}\mathbf{X}^{(2)}) = \begin{bmatrix} 33 & 36 \\ 36 & 48 \end{bmatrix}
}
$$

\item[(i)] $\text{Cov}(\mathbf{X}^{(1)}, \mathbf{X}^{(2)})$

\noindent This is simply the upper-right 2x2 block of the mean vector:

$$
\boxed{
\text{Cov}(\mathbf{X}^{(1)}, \mathbf{X}^{(2)}) = \boldsymbol{\Sigma}_{12} = \begin{bmatrix} 2 & 2 \\ 1 & 0 \end{bmatrix}
}
$$

\item[(j)] $\text{Cov}(\mathbf{A}\mathbf{X}^{(1)}, \mathbf{B}\mathbf{X}^{(2)})$

\vspace{0.5em}

\noindent Using $\text{Cov}(\mathbf{A}\mathbf{X}^{(1)}, \mathbf{B}\mathbf{X}^{(2)}) = \mathbf{A}\,\text{Cov}(\mathbf{X}^{(1)},\mathbf{X}^{(2)})\,\mathbf{B}'$:

Substituting the values we have 
$$
= \begin{bmatrix} 1 & 2 \end{bmatrix} \begin{bmatrix} 2 & 2 \\ 1 & 0 \end{bmatrix} \begin{bmatrix} 1 & 2 \\ -2 & -1 \end{bmatrix}
$$

Multiplying the first two matrices,

$$
\begin{bmatrix} 1 & 2 \end{bmatrix} \begin{bmatrix} 2 & 2 \\ 1 & 0 \end{bmatrix} =
\begin{bmatrix} (1)(2)+(2)(1) & (1)(2)+(2)(0) \end{bmatrix}
= \begin{bmatrix} 4 & 2 \end{bmatrix}
$$

Then multiplying the result with the last matrix,

$$
\begin{bmatrix} 4 & 2 \end{bmatrix}
\begin{bmatrix} 1 & 2 \\ -2 & -1 \end{bmatrix}
=
\begin{bmatrix} (4)(1)+(2)(-2) & (4)(2)+(2)(-1) \end{bmatrix}
=
\begin{bmatrix} 4-4 & 8-2 \end{bmatrix} 
= 
\begin{bmatrix} 0 & 6 \end{bmatrix}
$$

Therefore, 

$$
\boxed{
\text{Cov}(\mathbf{A}\mathbf{X}^{(1)}, \mathbf{B}\mathbf{X}^{(2)}) = \begin{bmatrix} 0 & 6 \end{bmatrix}
}
$$

\end{enumerate}

\vspace{2em}

\noindent\rule{\textwidth}{0.4pt}

\vspace{2em}

\noindent {\Large \textbf{Problem 2.41}} \hfill {\large \textbf{(5 points)}}

\vspace{1em}

\noindent You are given the random vector $\mathbf{X}' = [X_1, X_2, X_3, X_4]$ with mean vector $\boldsymbol{\mu}_X' = [3, 2, -2, 0]$ and variance-covariance matrix

$$
\boldsymbol{\Sigma}_X =
\begin{bmatrix}
3 & 0 & 0 & 0 \\
0 & 3 & 0 & 0 \\
0 & 0 & 3 & 0 \\
0 & 0 & 0 & 3
\end{bmatrix}
= 3\mathbf{I}
$$

\noindent Let

$$
\mathbf{A} =
\begin{bmatrix}
1 & -1 & 0 & 0 \\
1 & 1 & -2 & 0 \\
1 & 1 & 1 & -3
\end{bmatrix}
$$

\begin{enumerate}
\item[(a)] $E(\mathbf{A}\mathbf{X})$, the mean of $\mathbf{A}\mathbf{X}$

\vspace{0.5em}

\noindent Using $E(\mathbf{A}\mathbf{X}) = \mathbf{A}\boldsymbol{\mu}_X$, with $\boldsymbol{\mu}_X = \begin{bmatrix} 3 \\ 2 \\ -2 \\ 0 \end{bmatrix}$:

$$
\mathbf{A}\boldsymbol{\mu}_X =
\begin{bmatrix}
(1)(3)+(-1)(2)+(0)(-2)+(0)(0) \\
(1)(3)+(1)(2)+(-2)(-2)+(0)(0) \\
(1)(3)+(1)(2)+(1)(-2)+(-3)(0)
\end{bmatrix}
=
\begin{bmatrix}
3-2 \\
3+2+4 \\
3+2-2
\end{bmatrix}
$$

$$
\boxed{
E(\mathbf{A}\mathbf{X}) = \begin{bmatrix} 1 \\ 9 \\ 3 \end{bmatrix}
}
$$

\item[(b)] $\text{Cov}(\mathbf{A}\mathbf{X})$, the variances and covariances of $\mathbf{A}\mathbf{X}$

\vspace{0.5em}

\noindent Using $\text{Cov}(\mathbf{A}\mathbf{X}) = \mathbf{A}\boldsymbol{\Sigma}_X\mathbf{A}' = \mathbf{A}(3\mathbf{I})\mathbf{A}' = 3\mathbf{A}\mathbf{A}'$, we first compute $\mathbf{A}\mathbf{A}'$ using the rows of $\mathbf{A}$:

$$
\mathbf{r}_1 = \begin{bmatrix} 1 & -1 & 0 & 0 \end{bmatrix}, \quad
\mathbf{r}_2 = \begin{bmatrix} 1 & 1 & -2 & 0 \end{bmatrix}, \quad
\mathbf{r}_3 = \begin{bmatrix} 1 & 1 & 1 & -3 \end{bmatrix}
$$

$$
\mathbf{r}_1\mathbf{r}_1' = 1+1+0+0 = 2, \qquad
\mathbf{r}_1\mathbf{r}_2' = 1-1+0+0 = 0, \qquad
\mathbf{r}_1\mathbf{r}_3' = 1-1+0+0 = 0
$$

$$
\mathbf{r}_2\mathbf{r}_2' = 1+1+4+0 = 6, \qquad
\mathbf{r}_2\mathbf{r}_3' = 1+1-2+0 = 0, \qquad
\mathbf{r}_3\mathbf{r}_3' = 1+1+1+9 = 12
$$

$$
\mathbf{A}\mathbf{A}' = \begin{bmatrix} 2 & 0 & 0 \\ 0 & 6 & 0 \\ 0 & 0 & 12 \end{bmatrix}
$$

$$
\text{Cov}(\mathbf{A}\mathbf{X}) = 3\mathbf{A}\mathbf{A}' = 3\begin{bmatrix} 2 & 0 & 0 \\ 0 & 6 & 0 \\ 0 & 0 & 12 \end{bmatrix}
$$

$$
\boxed{
\text{Cov}(\mathbf{A}\mathbf{X}) = \begin{bmatrix} 6 & 0 & 0 \\ 0 & 18 & 0 \\ 0 & 0 & 36 \end{bmatrix}
}
$$

\item[(c)] Which pairs of linear combinations have zero covariance?

\vspace{0.5em}

\noindent Writing $\mathbf{Y} = \mathbf{A}\mathbf{X} = \begin{bmatrix} Y_1 \\ Y_2 \\ Y_3 \end{bmatrix}$, the off-diagonal entries of $\text{Cov}(\mathbf{A}\mathbf{X})$ found in part (b) give the pairwise covariances. Since all off-diagonal entries of $\text{Cov}(\mathbf{A}\mathbf{X})$ equal zero, i.e., $\text{Cov}(Y_1,Y_2) = \text{Cov}(Y_1,Y_3) = \text{Cov}(Y_2,Y_3) = 0$,

$$
\boxed{\text{every pair among } Y_1, Y_2, Y_3 \text{ has zero covariance.}}
$$

\end{enumerate}

\end{document}